\documentclass[14pt,usepdftitle=false,aspectratio=169]{beamer}
\usepackage{preambule}
\setbeamercolor{structure}{fg=black}
\usepackage{al}
\hypersetup{pdftitle={Surfaces de Riemann -- Monodromie}}
\title{Surfaces de Riemann\\\emph{Monodromie}}
\author{}
\date{}
\begin{document}
\begin{frame}
    \titlepage
\end{frame}
\slideq{Représentation de monodromie}{1/3}
\slider{Morphisme $\rho\!:\!\Pi_1\l\bbC\setminus S,y_0\r\to\gl{V_0}$ où $V_0$ est l'espace vectoriel des solutions holomorphes d'une équation différentielle homogène}{1/3}
\slideq{Représentation de monodromie d'un revêtement $f\!:\!X\to Y$ holomorphe}{2/3}
\slider{Application $\rho_f\!:\!\Pi_1\l Y\setminus B\l f\r,y_0\r\to\frakS_{\operatorname{deg}\l f\r}$ définie par $\rho_f\l\gamma\r\l i\r=j$ où si $f^{-1}\l y_0\r=\set{x_1,\cdots,x_d}$, alors le relèvement de $\gamma$ partant en $x_i$ se termine en $x_j$}{2/3}
\slideq{Prolongement analytique de $u$ le long d'un chemin $\gamma$}{3/3}
\slider{Famille d'applications holomorphes $\l u_t\!:\!U_t\to\bbC\r_{t\in\lc0,1\rc}$ avec $U_t$ un voisinage de $\gamma\l t\r$ et telle que, pour tout $t_0\in\lc0,1\rc$, il existe $\delta>0$ tel que, pour tout $t\in\left]t_0-\delta,t_0+\delta\right[$, $u_t$ et $u_{t_0}$ coïncident sur $U_t\cap U_{t_0}$}{3/3}
\end{document}